\section{Exponential confinement by physical anchor layers}
\label{sec:ae-exponential}
The polynomial estimates above do not exhaust the variational mechanism.
The physical turning bound supplies contacts at which the endpoint work
of a whole comparison interval can be paid by energy in its removed
boundary layer.  This yields a strict contraction after a fixed number
of actual collisions.  The resulting exponential estimate concerns
every global disk-product minimum, including nonphysical and inactive
ones.  It is not obtained by integrating a local endpoint derivative
along an assumed physical continuation.

Keep the original reference, comparison, radius and constants of
Sections~\ref{sec:bl-global} and~\ref{sec:tb-clearance}.  In particular
$Y$ is a global minimum of the unchanged disk-word problem, and the
reference is actual regular or a fixed-word closed physical reference.
To distinguish it from the depth of an index, denote the anchor threshold
by $a_0=1/416640$.  Thus $N=94$, $r=47/100$, $\tau_-=3/50$, and
\begin{equation}\label{eq:ae-contraction}
 \kappa=\frac{r}{a_0\tau_-}=3263680,
 \qquad q_{\rm a}=\frac{\kappa}{1+\kappa}
                 =\frac{3263680}{3263681}\in(0,1).
\end{equation}
The symbol $q_{\rm a}$ is a deterministic contraction constant; it
does not replace the original guard constant $q_0=47/53$.

\begin{lemma}[Common active geometry of global minima]
\label{lem:ae-common-minima}
For a fixed word and radius, all global disk-product minima have the
same ordered unit edge directions and the same active index set.
At every active index their vertices agree exactly.  Both endpoints
are active.  Any difference between minima lies along the common
straight direction in a chain of inactive contacts.  This is a
convexity fact about the comparison problem, not a new billiard map
or a claim that its full minimum is unique.
\end{lemma}
\begin{proof}
Let $Y,Z$ be two minima.  Their midpoint lies in the same product of
disks and, by convexity of $A$, is another minimum.  The sum of the
nonnegative defects in the flightwise triangle inequalities is zero,
so each such inequality is an equality.  Both vectors of each edge
are nonzero, by the fixed consecutive-disk separation.  Equality in
the Euclidean triangle inequality therefore makes them positive
parallel vectors, and their unit directions agree.

The Cartesian gradients at every coordinate consequently agree for
all minima.  If this common gradient $g_i$ is nonzero, the minimum
condition on its disk gives the unique point
$O_i-Rg_i/|g_i|$.  If it is zero, its neighboring directions agree;
this index is inactive for every minimum.  At an endpoint the gradient
is a unit vector up to sign, so it is always active.  On a consecutive
chain of inactive indices the edge direction is constant, and the
differences between corresponding edges of $Y$ and $Z$ are multiples
of that direction.  Starting at the preceding common active vertex
and summing these differences shows that all possible vertex
differences on the chain are longitudinal.  No continuation through
an inactive vertex or through an unlisted disk is inferred.
\end{proof}

\begin{lemma}[The exact work of an anchor-bounded interval]
\label{lem:ae-hard-interval}
If $1\le u<v\le m-1$ and both reference contacts are anchors,
$c_u,c_v\ge a_0$, then
\begin{equation}\label{eq:ae-hard-interval}
 \frac1R\sum_{i=u+1}^{v-1}H_i+
           \frac12\sum_{i=u}^{v-1}F_i
       \le\frac{H_u+H_v}{a_0\tau_-}.
\end{equation}
Both flights adjacent to the cut are retained on the left.  In
particular the inequality also covers $v=u+1$.
\end{lemma}
\begin{proof}
The variational comparison from Lemma~\ref{lem:bl-caccioppoli} gives
$(e_i-e_{i-1})\cdot\delta_i\ge H_i/R$ for $u<i<v$.
Summing on this hard interval, with its endpoint work retained, yields
\begin{equation}\label{eq:ae-hard-work}
 \frac1R\sum_{i=u+1}^{v-1}H_i+
 \frac12\sum_{i=u}^{v-1}(\ell_i+\tau_i)|e_i|^2
 \le e_{v-1}\cdot\delta_v-e_u\cdot\delta_u.
\end{equation}
Indeed the boundary of the finite integration by parts is exactly
the right side, and
$e_i\cdot(\delta_{i+1}-\delta_i)
 =(\ell_i+\tau_i)|e_i|^2/2$.  This identity is still valid when
there are no interior contacts.

At either cut use the elementary inequality
\[
 |e|\,|\delta|\le\frac{\tau_-}{4}|e|^2
                           +\frac{|\delta|^2}{\tau_-}.
\]
An edge incident to only one cut loses at most $\tau_-|e_i|^2/4$
from its kinetic coefficient; its remaining coefficient is at least
$\ell_i/2$.  If $v=u+1$, the same edge pays both cuts, but
\[
 \frac{\ell_i+\tau_i}{2}-\frac{\tau_-}{2}
                         \ge\frac{\ell_i}{2}.
\]
Thus every edge retains at least $F_i/2$, including the two boundary
flights.  At the actual reference anchors,
$|\delta_u|^2\le H_u/a_0$ and
$|\delta_v|^2\le H_v/a_0$.  These facts prove
\eqref{eq:ae-hard-interval}.  No assertion about the comparison
incidences is used.
\end{proof}

For an integer $d\ge0$ define the full relative energy at depth $d$ by
\begin{equation}\label{eq:ae-depth-energy}
 \mathcal E_d(Q,Y)=
 \frac1R\sum_{\substack{0<i<m\\d\le i\le m-d}}H_i+
                  \frac12\sum_{i=d}^{m-d-1}F_i.
\end{equation}
Every sum over an empty interval is zero.  In particular a central
single contact is kept when $m=2d$, although no flight remains.
The actual endpoints are never assigned interior contact weights.

\begin{theorem}[A contraction on every physical anchor layer]
\label{thm:ae-layer}
For every integer $d\ge0$,
\begin{equation}\label{eq:ae-layer}
 \mathcal E_{d+N}\le q_{\rm a}\mathcal E_d,
 \qquad
 \mathcal E_d\le E_R(Q)q_{\rm a}^{\lfloor d/N\rfloor}.
\end{equation}
The inequalities hold simultaneously for every global minimum and
every original word and radius, including at a closed physical
reference with zero internal incidences.  They have no collision-count
prefactor and no inverse individual incidence.
\end{theorem}
\begin{proof}
The complete positive gap, before any separate estimate of its
summands, gives
\begin{equation}\label{eq:ae-initial-energy}
 \mathcal E_0
 =\frac1R\sum_{i=1}^{m-1}H_i+\frac12\sum_{i=0}^{m-1}F_i
 \le\sum_{i=1}^{m-1}h_i+\frac12\sum_{i=0}^{m-1}F_i
 \le G\le E_R(Q).
\end{equation}
Thus both kinds of energy consume the same positive budget, not
two copies of it.  The last inequality uses global minimality only
through $Y\in\mathcal S(Q)$.

If $2(d+N)>m$, the first inequality in \eqref{eq:ae-layer} has zero
left side.  Otherwise the two complete physical $N$-flight blocks
starting at $d$ and $m-d-N$ supply reference anchors
\[
 u\in\{d+1,\ldots,d+N-1\},\qquad
 v\in\{m-d-N+1,\ldots,m-d-1\}.
\]
These are interior contacts, and $u<d+N\le m-d-N<v$.
Consequently every energy summand in $\mathcal E_{d+N}$ occurs
on the left of \eqref{eq:ae-hard-interval}.  Conversely, the two
anchor summands $H_u/R,H_v/R$ occur in
$\mathcal E_d-\mathcal E_{d+N}$.  All the other summands in that
difference are nonnegative.  Lemma~\ref{lem:ae-hard-interval} gives
\[
 \mathcal E_{d+N}
 \le\frac{H_u+H_v}{a_0\tau_-}
 =\frac{R}{a_0\tau_-}\frac{H_u+H_v}{R}
 \le\kappa(\mathcal E_d-\mathcal E_{d+N}).
\]
Rearranging proves the strict contraction with exactly
$q_{\rm a}=\kappa/(1+\kappa)$.  The argument also covers
$m=2(d+N)$, where the inner energy contains one contact.  No boundary
flight is counted as a removed-layer payment: its Young payment has
already been absorbed in \eqref{eq:ae-hard-interval}.

Iterate the contraction at the depths $0,N,2N,\ldots$ and use
\eqref{eq:ae-initial-energy}.  Since $\mathcal E_d$ is nonincreasing,
$\mathcal E_d\le\mathcal E_{N\lfloor d/N\rfloor}$.
This proves the second inequality.  The anchors can be chosen anew
at each layer; no compatible continuation or fixed selection of
anchor points is required.
\end{proof}

Put
\begin{equation}\label{eq:ae-j-and-B}
 j(d)=\max\{0,\lfloor d/N\rfloor-1\},
 \qquad B_\perp=B_G+B_E.
\end{equation}

\begin{corollary}[Exponential position, direction and transverse budgets]
\label{cor:ae-geometric}
For every integer $d\ge0$ and every global minimum,
\begin{align}
 \sum_{i=d}^{m-d}|\delta_i|^4
     &\le C E_R(Q)q_{\rm a}^{j(d)},
                       \label{eq:ae-position}\\
 \sum_{i=d}^{m-d-1}F_i
     &\le 2E_R(Q)q_{\rm a}^{\lfloor d/N\rfloor},
                       \label{eq:ae-flight}\\
 \sum_{i=d}^{m-d-1}\beta_i^2
     &\le B_\perp E_R(Q)q_{\rm a}^{j(d)}.
                       \label{eq:ae-transverse}
\end{align}
The fourth-power and transverse sums have their usual zero values on
empty intervals.  The position estimate includes both actual endpoints
when $d=0$; they are paid by the original deficit.
\end{corollary}
\begin{proof}
The flight estimate follows directly from the definition of
$\mathcal E_d$.  For $d<N$, the position and transverse estimates
are respectively Theorem~\ref{thm:ss-quartic} and
\eqref{eq:tb-global-budget} on the sublevel.

For $d\ge N$ with a nonempty target interval, apply the interior
coercivity lemma on $[u,v]=[d-N,m-d+N]$.  Its target range contains
all the positions in \eqref{eq:ae-position}.  The positive contact
and flight components are jointly bounded by $\mathcal E_{d-N}$,
so
\[
 S_d\le\max\{rA_H,2A_F\}\mathcal E_{d-N}
       \le C\mathcal E_{d-N}.
\]
For the constants in \eqref{eq:bl-local-constants}, the last comparison
is the explicit rational inequality
\[
 \max\{rA_H,2A_F\}<10^8=C.
\]
Similarly the local transverse
estimate \eqref{eq:tb-local-budget} contains all the included flights
and gives
\[
 \sum_{i=d}^{m-d-1}\beta_i^2
 \le\max\{rB_H,2B_F\}\mathcal E_{d-N}
 \le B_\perp\mathcal E_{d-N},
\]
where $\max\{rB_H,2B_F\}<B_G+B_E$ is another direct rational
comparison of the displayed constants.  Both inequalities use the
shared energy before estimating its components.  Since
$\lfloor(d-N)/N\rfloor=j(d)$ for $d\ge N$, the layer theorem
proves the two conclusions.  The polynomial bounds already proved
remain valid and may be used together with these exponential ones;
no preceding proof body is removed.
\end{proof}

\begin{corollary}[Explicit exponentially small physical failure witnesses]
\label{cor:ae-witnesses}
For an interior contact of depth $d_i\ge1$, inactivity of a global
minimum implies
\begin{equation}\label{eq:ae-inactive-witness}
 c_i(Q)^2\le\frac{E_R(Q)}{\tau_-}
           q_{\rm a}^{\lfloor(d_i-1)/N\rfloor}.
\end{equation}
At every flight,
\begin{equation}\label{eq:ae-clearance-witness}
 \beta_i^2\le B_\perp E_R(Q)q_{\rm a}^{j(d'_i)}.
\end{equation}
If the last upper bound is strictly below $\tau_-^2/4$, a third-disk
violation of that comparison flight necessarily implies
$\Gamma_i(Q)\le\sqrt{B_\perp E_R(Q)}q_{\rm a}^{j(d'_i)/2}$.
All these witnesses concern the true incidence and clearance of the
original reference.  No probability or roof-density assertion follows
without a separate source argument.
\end{corollary}
\begin{proof}
Inactivity gives $2c_i\le|e_i|+|e_{i-1}|$, as before.  Both flights
belong to the energy at depth $d_i-1$, so
\[
 c_i^2\le\frac{|e_i|^2+|e_{i-1}|^2}{2}
       \le\frac1{2\tau_-}\sum_{l=d_i-1}^{m-d_i}F_l
       \le\frac{E_R(Q)}{\tau_-}
                    q_{\rm a}^{\lfloor(d_i-1)/N\rfloor}.
\]
The transverse point bound is a single summand of
\eqref{eq:ae-transverse}.  Its stated strict regime permits application
of the already proved strip lemma, giving the clearance implication.
This excludes no original collision, and treats every minimizer with
the same constants.
\end{proof}

\begin{remark}[Rate, original guard, and probability topology]
The contraction $q_{\rm a}$ is close to one because the explicit
physical anchor bound is conservative.  It is nevertheless fixed
before any collision limit, and its exponential decay is uniform on
the entire original geometric family.  It does not identify these
rates with the much narrower original guard widths
$\varepsilon q_0^{d/4}$.  In particular a substitution of one depth
schedule for the other would change the claimed microscopic problem
and is not made.  Exponential geometric cost, exact original-source
mass, inherited-pressure conditional roof height, and the entire
normalized shortage remain distinct assertions.  None is inferred
from a finite trajectory computation.
\end{remark}
